\bookchapter{Приложение H: дорожная карта версий}

\chapterintro{Эта книга постоянно обновляется. Ниже — планы на будущие версии. Будем рады, если в отзывах вы расскажете, о чём вам хочется прочитать в первую очередь.}

\textbf{Реализовано в v1.8: добавлены части VII, VIII и IX — «Программирование с ИИ на практике», «Разработка под Android на практике» и «Разработка под iOS на практике» (главы 37–69, всего 33 главы)}\par

От парадигмальной революции в программировании с ИИ, пяти принципов мышления, инженерии контекста и требований, панорамы инструментов и моделей 2026 года, практики работы с API, инженерии диалога, трёх полноценных проектов, обеспечения качества и пути к мастерству — вы пройдёте путь от «умею пользоваться ИИ» до «руковожу ИИ, который пишет код»; а далее — полный цикл разработки на двух платформах, Android и iOS: настройка окружения (Android Studio / Xcode), экспресс-курс языков (Kotlin / Swift), фреймворки интерфейсов (Jetpack Compose / SwiftUI), данные и сеть, практические проекты (учёт расходов + приложение погоды), тестирование и публикация (Google Play / App Store) и справочники — путь от цифрового пользователя до разработчика приложений.

\begin{tcbinfo}{Планы на следующую версию (v1.9)}
\begin{itemize}
  \item Смартфон и носимые устройства: гид по глубокой связке с Apple Watch, часами Huawei и умными кольцами (Oura / Galaxy Ring)
  \item Путешествия и смартфон: практическая работа с eSIM, варианты международного интернета, офлайн-карты, оплата за рубежом, сравнение приложений-переводчиков
  \item Смартфон и здоровье: спорт и активность, мониторинг сна, пульс и кислород в крови, женское здоровье, приложения для ментального здоровья
  \item Создание коротких видео с ИИ: практические уроки по CapCut AI, Kling AI, цифровым аватарам и ИИ-озвучке
  \item Введение в вертикальные мини-сериалы: структура сценария, техника съёмки, правила платформ
\end{itemize}

\end{tcbinfo}

\begin{tcbwarn}{Примечание об актуальности}
Технологии смартфонов развиваются очень быстро, и книга написана по состоянию технологий и рынка на середину 2026 года. Возможности ИИ, модели чипов, версии систем и ценовые диапазоны к моменту вашего чтения могли уже измениться. В конкретных шагах ориентируйтесь на актуальную версию системы вашего устройства, а при принятии решений о покупке сверяйтесь с последними обзорами и отзывами пользователей.

\end{tcbwarn}
